\clearpage
\section{Semantic weakening with retained free resources}
\label{sec:resource-semantic-weakening}
This milestone removes the restriction on other free resources in the boxed
weakening calculation. It also proves the unary-operation and abstraction
induction steps for the original Figure~16 translation, for both variable
and atom extension. No term, translation clause, canonical distributor,
or machine rule has changed.

The connection to the paper is the proof of the weakening laws in
Section~12.3 (HTML Proof~12.3 of Proposition~11; PDF Proposition~12.1).
Its third local situation pulls an unused output through a box and then
reduces its unary distributor using Idem. The present construction retains
all other free-resource rows while performing exactly that calculation.
The original source is
\href{https://arxiv.org/html/1907.01257v5\#S12.SS3}{Section~12.3 of the paper}.

\paragraph{Scope of the result.}
The completed box law permits every binder count $n$, every old variable
count $k$, every old atom count $h$, and every depth allowance $d$.
The old body may contain arbitrary nested boxes and may use its old free
resources. The proof supplies concrete filling certificates and hence exact
contextual equivalence; it assumes neither successful termination nor
robustness of a new template. The base translation cases and assembly of the complete term induction
remain open. Section~\ref{sec:binary-semantic-weakening} subsequently
completes extraction of the selected binary row for independent subterms. Thus this is a completed part of
semantic weakening, not yet Proposition~12.2.

\Needspace{12\baselineskip}
\subsection{Serial composition of ordered parallel families}
\termdef{serial-parallel-interchange}
Let $F_a:I_a\to J_a$ and $R_a:J_a\to O_a$ be finite families of hypernets.
A prefix boundary $C$ orders all the $I_a$ and $J_a$ ports; a suffix boundary
$B$ orders all the $J_a$ and $O_a$ ports. Assume that $C$'s output bijection
is exactly $B$'s input bijection. The composite boundary uses $C$'s input
bijection and $B$'s output bijection. The two presentations are
\[
 \Bigl(\mathop{\parallel}_C F_a\Bigr);
 \Bigl(\mathop{\parallel}_B R_a\Bigr)
 \quad\hbox{and}\quad
 \mathop{\parallel}_{C,B}(F_a;R_a).
\]
Unlike the earlier input-free interchange, the $I_a$ are arbitrary here.
\lean{Copying/ParallelSerial}{16}{20}{ParallelSerial.afterBoundary}
\lean{Copying/ParallelSerial}{22}{22}{ParallelSerial.left}
\lean{Copying/ParallelSerial}{23}{23}{ParallelSerial.right}

\begin{theorem}[Ordered serial-parallel interchange]
These two presentations are recursively isomorphic with every external
boundary position fixed.
\end{theorem}
\begin{proof}
Send a shared global wire $j$ to the unique component and local port selected
by $B$'s inverse input bijection. Reassociate a prefix internal vertex
$(\mathsf{left},a,v)$ as $(a,\mathsf{left},v)$, and likewise for suffix
vertices and all edges. The inverse maps use $B$'s forward input bijection.
Its two inverse laws prove that these maps are bijections. The shared-wire
label law is precisely $B$'s input typing law. For incidence, prefix outputs
and suffix inputs become the same shared internal vertex; the compatibility
assumption identifies their two original orderings. Original external inputs
and outputs remain external at the same positions. Mapping each source and
target list through these commuting port maps proves the isomorphism.
\end{proof}
\lean{Copying/ParallelSerial}{131}{161}{ParallelSerial.iso}

\Needspace{12\baselineskip}
\subsection{Retaining resources in the original distributors}
\termdef{resource-distributor-embedding}
Write $E_{k,h}=*^k\diamond^h$. Adding a resource of sort $s$ means adding a
variable at the end of the variable block, or adding an atom at the start
of the atom block. In either case the new global output is at position $k$.
Old outputs before $k$ stay fixed; old outputs at or after $k$ move one
position to the right. This is the existing BPullW insertion convention.

For every number of uses $m$, a canonical term row has its literal occurrence
in the term half of $D_{k,m}^{*}\otimes D_{h,m}^{\diamond}$; an atom row has
its literal occurrence in the atom half. The maps tag each internal vertex
and edge by its sort and resource index. Their incidence laws are the
original parallel embeddings. The selected new row is the last term row
or the first atom row of the enlarged distributor.
\lean{Copying/DistributorExtension}{11}{31}{DistributorExtension.termOccurrence}
\lean{Copying/DistributorExtension}{34}{54}{DistributorExtension.atomOccurrence}
\lean{Copying/DistributorExtension}{58}{63}{DistributorExtension.selected}

Assemble all the old term rows into one occurrence using their original
ordered distributor boundary, and do the same for the old atom rows.
Their resource injections are injective, so distinct rows have disjoint
edges. Their term/atom tags are different, so the two complete forests are
also disjoint. Their parallel assembly is the occurrence of the entire
old environment distributor in the enlarged distributor.
\lean{Copying/DistributorExtension}{65}{74}{DistributorExtension.termForest}
\lean{Copying/DistributorExtension}{76}{85}{DistributorExtension.atomForest}
\lean{Copying/DistributorExtension}{87}{99}{DistributorExtension.distributor}

The four component-port lemmas say that following a port first into its
old forest and then into the enlarged distributor gives exactly the same
port as entering its corresponding new row directly. They follow from the
parallel-assembly port equation, first within each sort and then across
the term/atom split. They apply to internal, input and output ports.
\lean{Copying/DistributorExtension}{101}{110}{DistributorExtension.termForest_port}
\lean{Copying/DistributorExtension}{112}{121}{DistributorExtension.atomForest_port}
\lean{Copying/DistributorExtension}{123}{130}{DistributorExtension.distributor_term_port}
\lean{Copying/DistributorExtension}{132}{139}{DistributorExtension.distributor_atom_port}

Every old output follows the insertion map just described. More precisely,
the input for old variable $i$ at use $u$ originally has coordinate $i+ku$;
after extension it has coordinate $\iota_s(i)+k'u$. For atoms the analogous
coordinate is $j+hu$ inside the atom block, with that block following all
$k'm$ term inputs in the enlarged graph. These are the original use-major
coordinates, not a resource-major reordering. The output theorem splits at
$k$; the two input theorems instantiate the component-port lemmas at use $u$.
\lean{Copying/DistributorExtension}{141}{155}{DistributorExtension.distributor_output}
\lean{Copying/DistributorExtension}{157}{162}{DistributorExtension.distributor_term_input}
\lean{Copying/DistributorExtension}{164}{172}{DistributorExtension.distributor_atom_input}

\Needspace{12\baselineskip}
\subsection{Exhaustive extraction of the added unary row}
\termdef{selected-unused-row}
Let $R_{k,h}$ denote the original single-use environment distributor. Its
input equalities $k\cdot1=k$ and $h\cdot1=h$ can be expressed by explicit
boundary transport. The auxiliary raw presentation changes only this
transport; the original single-use graph equals it.
\lean{Copying/SingleUseExtension}{12}{13}{SingleUseExtension.raw}
\lean{Copying/SingleUseExtension}{15}{16}{SingleUseExtension.original}

Specializing the resource-preserving embedding to $m=1$ and transporting
its boundaries gives the retained occurrence. Both its inputs and outputs
follow the same old-resource insertion map. For inputs, split the interface
into its term and atom blocks and specialize the use coordinate to zero;
for outputs, use the general output theorem.
\lean{Copying/SingleUseExtension}{19}{21}{SingleUseExtension.retained}
\lean{Copying/SingleUseExtension}{23}{28}{SingleUseExtension.retained_output}
\lean{Copying/SingleUseExtension}{30}{54}{SingleUseExtension.retained_input}

The added occurrence is the selected original $D_{1,1}^{s}$ row. Its sole
input and output both map to the new global position $k$. The boundary for
assembling the retained forest beside this row uses the same typed insertion
bijection on inputs and outputs. Its old-port and added-port equations
therefore agree exactly with those of the two occurrences.
\lean{Copying/SingleUseExtension}{57}{59}{SingleUseExtension.added}
\lean{Copying/SingleUseExtension}{61}{70}{SingleUseExtension.added_input}
\lean{Copying/SingleUseExtension}{72}{81}{SingleUseExtension.added_output}
\lean{Copying/SingleUseExtension}{95}{108}{SingleUseExtension.boundary}
\lean{Copying/SingleUseExtension}{111}{113}{SingleUseExtension.boundary_old}
\lean{Copying/SingleUseExtension}{115}{119}{SingleUseExtension.boundary_added}

The component family consists of the retained occurrence and the added row.
They are disjoint: an old variable index is strictly below the new last
variable, while an old atom index after insertion is positive and the new
atom index is zero. Opposite-sort edges have different tags. Parallel
assembly gives their combined occurrence; the boundary equations prove
that it fixes every external input and output.
\lean{Copying/SingleUseExtension}{122}{125}{SingleUseExtension.components}
\lean{Copying/SingleUseExtension}{143}{154}{SingleUseExtension.occurrence}

\begin{lemma}[Exhaustive row coverage]
Every edge of the original enlarged single-use distributor belongs to one
of these two occurrences.
\end{lemma}
\begin{proof}
For variable extension, split a variable row index into $i<k$ and $i=k$;
all atom rows are retained. For atom extension, split an atom row index
into $i=0$ and $i>0$; in the latter case its old index is $i-1$. All variable
rows are retained. Each case preserves the original within-row edge index.
This includes empty old variable or atom blocks. Complete incidence then
supplies internal-vertex coverage as well. The existing occurrence-cover
lemma upgrades this edge-surjective, boundary-fixing occurrence to an
ordered recursive isomorphism.
\end{proof}
\lean{Copying/SingleUseExtension}{182}{218}{SingleUseExtension.edge_coverage}
\lean{Copying/SingleUseExtension}{222}{229}{SingleUseExtension.iso}

\Needspace{12\baselineskip}
\subsection{The complete single-use weakening calculation}
Let $W_sG$ denote the original BPullW extension of $G$ by a separate weakening.
For an arbitrary $G:I\to E_{k,h}$, the prefix assembly places $G$ beside
that weakening, retaining $I$ unchanged and ordering its outputs by the
same insertion bijection. The extension isomorphism identifies this
presentation with $W_sG$ by identity maps on edges and internal vertices;
only the explicitly equal boundary presentations differ.
\lean{Copying/WeakeningSingleUse}{13}{24}{Weakening.SingleUse.prefixBoundary}
\lean{Copying/WeakeningSingleUse}{27}{48}{Weakening.SingleUse.extension_iso}

\begin{theorem}[Single-use extension law]
If $G$ is focus-free, there is a genuine Idem filling witnessing
\[
       (W_sG);R_{k',h'}\ \leadsto\ W_s(G;R_{k,h}),
\]
where $(k',h')$ is the enlarged environment. In particular, the two graphs
are connected by a weakening calculation.
\end{theorem}
\begin{proof}
Extract the added row from $R_{k',h'}$, then apply serial-parallel
interchange. The resulting components are the unchanged $G;R_{k,h}$ and
one incoming weakening followed by the actual $D_{1,1}^{s}$. The latter
still contains its own canonical weakening and contraction. Apply the
proved unary Idem witness to this component, and a reflexive witness to
the old component. Parallel assembly and the endpoint isomorphisms give
one original whole-graph Idem filling. Its calculation is simply the Idem
constructor applied to that certificate.
\end{proof}
\lean{Copying/WeakeningSingleUse}{51}{65}{Weakening.SingleUse.feeding_iso}
\lean{Copying/WeakeningSingleUse}{69}{85}{Weakening.SingleUse.filling}
\lean{Copying/WeakeningSingleUse}{89}{93}{Weakening.SingleUse.normalize}

\Needspace{12\baselineskip}
\subsection{Boxes with arbitrary other free resources}
\termdef{boxed-resource-weakening}
\begin{theorem}[General boxed weakening]
Let $G:*\to E_{n+k,h}$ be any focus-free body, with its first $n$ term
outputs bound. Let $B(G):\mathsf{thunk}(n)\to E_{k,h}$ be its original box.
Then the following is an exact weakening calculation:
\[
 B(W_sG);R_{k',h'}\ \leadsto\ W_s\bigl(B(G);R_{k,h}\bigr).
\]
Here $W_sG$ denotes the existing bound-aware BPullW body constructor, so
none of the first $n$ bound outputs is moved.
\end{theorem}
\begin{proof}
Use the literal BPullW certificate to move the weakening outside the box,
retaining the outer distributor as an unchanged serial suffix. Apply the
single-use extension law to the original box. The result retains the whole
body $G$ and all its old free-resource rows. The two stages are therefore
one BPullW filling followed by one Idem filling, with only ordered graph
isomorphisms used to expose the selected row.
\end{proof}
\lean{Copying/WeakeningTranslationSteps}{12}{21}{Weakening.boxed_resources}

The addressed theorem plugs this complete calculation into any recursive
site of a focus-free closing host. Hole freedom of the filled source is the
only hole-freedom premise; the calculation supplies it for the target and
intermediates. Exact semantic soundness then proves both refinements with
equality of successful-transition counts in the original unbounded machine.
No binding-free assumption is needed for this structural result; it therefore
also applies to all the intended binding-free observation contexts.
\lean{Copying/WeakeningTranslationSteps}{89}{101}{Weakening.boxed_resources_addressed}

\Needspace{12\baselineskip}
\subsection{The original unary and abstraction induction steps}
\termdef{unary-abstraction-weakening}
For a unary operation $o$, write $U_o(G)$ for the unchanged Figure~16
clause: the operation edge, then $G$, then $R_{k,h}$. The extension law
\[
 U_o(W_sG)\leadsto W_s(U_o(G))
\]
first moves the separate weakening past the unchanged operation prefix by
serial-extension isomorphism, then applies the single-use extension law.
It covers negation, allocation and dereference without a special operational
assumption. Given an actual body calculation $G'\leadsto W_sG$, serial
lifting through the operation and enlarged distributor, followed by this
law, gives $U_o(G')\leadsto W_s(U_o(G))$.
\lean{Copying/WeakeningTranslationSteps}{25}{36}{Weakening.unary_extend}
\lean{Copying/WeakeningTranslationSteps}{40}{49}{Weakening.unary_step}

For abstraction, the original body has $k+1$ term outputs. The existing
\texttt{betaBody} constructor expresses unused-resource insertion in that
convention while preserving the distinguished bound output. The abstraction
extension law uses the proved BPullW lambda-payload witness, keeps the outer
single-use distributor, and then applies the same row correction. This
proves $A(W_sG)\leadsto W_s(A(G))$ for the original abstraction clause $A$.
\lean{Copying/WeakeningTranslationSteps}{53}{61}{Weakening.abstraction_extend}

Calculation congruence for abstraction first transports the equality
$k+1=1+k$, boxes the body calculation, adds the unchanged lambda head,
and appends the unchanged single-use distributor. This is a lift of the
actual proof tree; no equivalence hypothesis is introduced. Consequently,
from a body calculation $G'\leadsto\mathtt{betaBody}(G,s)$, abstraction
congruence followed by the extension law gives
$A(G')\leadsto W_s(A(G))$. Focus freedom of $G$ is recovered from the
supplied calculation and the proved extension invariant.
\lean{Copying/WeakeningTranslationSteps}{65}{75}{Weakening.Calculation.abstraction}
\lean{Copying/WeakeningTranslationSteps}{79}{85}{Weakening.abstraction_step}

\Needspace{12\baselineskip}
\subsection{Checked instances and remaining dependency}
The regression body is the direct translation of $\lambda z.\,(v_1\ a_1)$
in an environment with two variables and two atoms: under its binder the
variable index is $2$ and the atom index is $1$. It already contains a box
and uses both an outer variable and an ambient atom. Its recursive focus
freedom follows from the general translation theorem.
\lean{Copying/WeakeningTranslationExamples}{11}{12}{Weakening.Examples.nestedBody}
\lean{Copying/WeakeningTranslationExamples}{14}{15}{Weakening.Examples.nestedBody_free}

The mixed-box instance encloses this graph in another box while keeping one
old free variable and both atoms. The mixed-abstraction instance uses the
actual lambda translation clause. Both are checked for either added sort.
The allocation/read instance composes the allocation extension law with
the dereference induction step. These are structural semantic calculations,
not claims that every supplied program terminates.
\lean{Copying/WeakeningTranslationExamples}{19}{25}{Weakening.Examples.mixed_box}
\lean{Copying/WeakeningTranslationExamples}{29}{32}{Weakening.Examples.mixed_abstraction}
\lean{Copying/WeakeningTranslationExamples}{36}{39}{Weakening.Examples.allocation_read}

The finite boundary checks verify that the new resource is at position $2$
when $k=h=2$. The second old atom moves from position $3$ to position $4$;
for atom extension, the old variable at position $1$ stays fixed. Decidable
finite computation proves these equations. In particular, the atom case
would fail if the new atom were incorrectly appended after the old atoms.
\lean{Copying/WeakeningTranslationExamples}{43}{45}{Weakening.Examples.term_boundary_order}
\lean{Copying/WeakeningTranslationExamples}{49}{52}{Weakening.Examples.atom_boundary_order}

\paragraph{Dependency map.}
The proved path is: original distributor occurrences $\to$ disjoint retained
and added rows $\to$ exhaustive unary-forest isomorphism $\to$ ordered
serial-parallel interchange $\to$ genuine whole-graph Idem filling
$\to$ BPullW plus Idem for general boxes $\to$ unary and lambda induction
steps $\to$ exact addressed observations. The semantic leaves are the
already-proved original Idem and BPullW templates. The input focus-freedom
conditions are application conditions, not missing robustness proofs.

What is still absent is the whole derivation for every translated term:
the variable, atom and constant base cases and the full syntactic induction
must be assembled before arbitrary-environment weakening can be claimed.
The subsequent Section~\ref{sec:binary-semantic-weakening} proves the
independent application/binary steps and their padding compatibility. Figure~22 and the final general-environment Proposition~12.2
remain downstream, together with the correspondence audits in Appendix~B.
